Information:

Picture 1 presents a negative image of sky taken by a CCD camera attached to
a telescope whose parameters are presented in the accompanying table (which
is part of the FITS datafile header). Picture 2 consists of two images: one
is an enlarged view of part of Picture 1 and the second is an enlarged image
of the same part of the sky taken some time earlier. Picture 3 presents a
sky map which includes the region shown in the CCD images.

The stars in the images are far away and should ideally be seen as point
sources. However, diffraction on the telescope aperture and the effects of
atmospheric turbulence (known as `seeing') blur the light from the stars.
The brighter the star, the more of the spread-out light is visible above the
level of the background sky.

Questions:

\begin{enumerate}
\item Identify any 5 bright stars (mark them by Roman numerals) from the
  image and mark them on both the image and map.
\item Mark the field of view of the camera on the map.
\item Use this information to obtain the physical dimensions of the CCD chip
  in mm.
\item Estimate the size of the blurring effect in arcseconds by examining
  the image of the star in Picture 2. (Note that due to changes in contrast
  necessary for printing, the diameter of the image appears to be about 3.5
  times the full width at half maximum (FWHM) of the profile of the star.)
\item Compare the result with theoretical size of the diffraction disc of
  the telescope.
\item Seeing of 1 arcsecond is often considered to indicate good conditions.
  Calculate the size of the star image in pixels if the atmospheric seeing
  was 1 arcsecond and compare it with the result from question 4.
\item Two objects observed moving relative to the background stars have been
  marked on Picture 1. The motion of one (``Object 1'') was fast enough that
  it left a clear trail on the image. The motion of the other (``Object 2'')
  is more easily seen on the enlarged image (Picture 2A) and another image
  taken some time later (Picture 2B).

  Using the results of the first section, determine the angular velocity on
  the sky of both objects. Choose which of the statements in the list below
  are correct, assuming that the objects are moving on circular orbits.
  (Points will be given for each correct answer marked and deducted for each
  incorrect answer marked.) The probable causes of the different angular
  velocities are:
  \begin{description}
  \item[a)] different masses of the objects,
  \item[b)] different distances of the objects from Earth,
  \item[c)] different orbital velocities of the objects,
  \item[d)] different projections of the objects' velocities,
  \item[e)] Object 1 orbits the Earth while Object 2 orbits the Sun.
  \end{description}
\end{enumerate}

Data: for Picture 1, the data are,

\begin{center}
\begin{tabular}{lll}
BITPIX & 16 & Number of bits per pixel \\
NAXIS & 2 & Number of axes \\
NAXIS1 & 1024 & Width of image (in pixels) \\
NAXIS2 & 1024 & Height of image (in pixels) \\
DATE-OBS & 2010-09-07 05:00:40.4 & Middle of exposure \\
TIMESYS & UT & Time Scale \\
EXPTIME & 300.00 & Exposure time (seconds) \\
OBJCTRA & 22 29 20.031 & RA of center of the image \\
OBJCTDEC & +07 20 00.793 & DEC of center of the image \\
FOCALLEN & 3.180m & Focal length of the telescope \\
TELESCOP & 0.61m & Telescope aperture \\
\end{tabular}
\end{center}

\ioaafig{2010_DA_D1-f1.pdf}

For Picture 2A (the same area observed some time earlier) the data are:
DATE-OBS = 2010-09-07 04:42:33.3 (middle of exposure). Picture 2B is an
enlargement of Picture 1 around Object 2.

\ioaafig{2010_DA_D1-f2.pdf}

\ioaafig{2010_DA_D1-f3.pdf}
